\documentclass[10pt,a4paper]{amsart}
\usepackage{amsmath,amssymb,amsthm,graphicx}
% Retained preamble for the Stacks Project chapter snapshot in this directory.
% The chapter's own preamble.tex is a copy of the shared Stacks Project preamble; the block
% below is sliced MECHANICALLY, line for line, from it. Nothing was added, removed or
% redefined; the wrapper class, the xr-hyper/hyperref wiring and the externaldocument list
% are omitted because the assembler supplies the document class.
%
% lines taken from preamble.tex: 3, 5, 9, 11, 12, 156, 157, 158, ... (51 lines total)
%
% package coverage restored from the chapter's own preamble (2026-10-13): xy, xy-2cell

\usepackage[all]{xy}
\xyoption{2cell}
\UseAllTwocells

\newenvironment{reference}{\comment}{\endcomment}
\newenvironment{slogan}{\comment}{\endcomment}
\newenvironment{history}{\comment}{\endcomment}
\theoremstyle{plain}
\newtheorem{theorem}[subsection]{Theorem}
\newtheorem{proposition}[subsection]{Proposition}
\newtheorem{lemma}[subsection]{Lemma}
\theoremstyle{definition}
\newtheorem{definition}[subsection]{Definition}
\newtheorem{example}[subsection]{Example}
\newtheorem{exercise}[subsection]{Exercise}
\newtheorem{situation}[subsection]{Situation}
\theoremstyle{remark}
\newtheorem{remark}[subsection]{Remark}
\newtheorem{remarks}[subsection]{Remarks}
\numberwithin{equation}{subsection}
\def\lim{\mathop{\mathrm{lim}}\nolimits}
\def\colim{\mathop{\mathrm{colim}}\nolimits}
\def\Spec{\mathop{\mathrm{Spec}}}
\def\Hom{\mathop{\mathrm{Hom}}\nolimits}
\def\Ext{\mathop{\mathrm{Ext}}\nolimits}
\def\SheafHom{\mathop{\mathcal{H}\!\mathit{om}}\nolimits}
\def\SheafExt{\mathop{\mathcal{E}\!\mathit{xt}}\nolimits}
\def\Sch{\mathit{Sch}}
\def\Mor{\mathop{\mathrm{Mor}}\nolimits}
\def\Ob{\mathop{\mathrm{Ob}}\nolimits}
\def\Sh{\mathop{\mathit{Sh}}\nolimits}
\def\NL{\mathop{N\!L}\nolimits}
\def\CH{\mathop{\mathrm{CH}}\nolimits}
\def\proetale{{pro\text{-}\acute{e}tale}}
\def\etale{{\acute{e}tale}}
\def\QCoh{\mathit{QCoh}}
\def\Ker{\mathop{\mathrm{Ker}}}
\def\Im{\mathop{\mathrm{Im}}}
\def\Coker{\mathop{\mathrm{Coker}}}
\def\Coim{\mathop{\mathrm{Coim}}}
\DeclareMathSymbol{\boxtimes}{\mathbin}{AMSa}{"02}
\def\QCohstack{\mathcal{QC}\!\mathit{oh}}
\def\Cohstack{\mathcal{C}\!\mathit{oh}}
\def\Spacesstack{\mathcal{S}\!\mathit{paces}}
\def\Quotfunctor{\mathrm{Quot}}
\def\Hilbfunctor{\mathrm{Hilb}}
\def\Curvesstack{\mathcal{C}\!\mathit{urves}}
\def\Polarizedstack{\mathcal{P}\!\mathit{olarized}}
\def\Complexesstack{\mathcal{C}\!\mathit{omplexes}}
\def\Pic{\mathop{\mathrm{Pic}}\nolimits}
\def\Picardstack{\mathcal{P}\!\mathit{ic}}
\def\Picardfunctor{\mathrm{Pic}}
\def\Deformationcategory{\mathcal{D}\!\mathit{ef}}
\usepackage{cleveref}
\crefname{originallemma}{original lemma}{original lemmas}
\Crefname{originallemma}{Original lemma}{Original lemmas}
\makeatletter
\newlabel{coherent-lemma-vanshing-gives-ample}{{0B5U}{source}{}{sourcepointer-1}{}}
\newlabel{coherent-lemma-vanshing-gives-ample@cref}{{[originallemma][0][]0B5U}{[1][0][]source}{}{}{}}
\makeatother
\crefname{originallemma}{original lemma}{original lemmas}
\Crefname{originallemma}{Original lemma}{Original lemmas}
\makeatletter
\newlabel{varieties-lemma-general-degree-g-line-bundle}{{0B8Z}{source}{}{sourcepointer-2}{}}
\newlabel{varieties-lemma-general-degree-g-line-bundle@cref}{{[originallemma][0][]0B8Z}{[1][0][]source}{}{}{}}
\makeatother
\crefname{originallemma}{original lemma}{original lemmas}
\Crefname{originallemma}{Original lemma}{Original lemmas}
\makeatletter
\newlabel{varieties-lemma-regular-point-on-curve}{{0B8Y}{source}{}{sourcepointer-3}{}}
\newlabel{varieties-lemma-regular-point-on-curve@cref}{{[originallemma][0][]0B8Y}{[1][0][]source}{}{}{}}
\makeatother
\crefname{originallemma}{original lemma}{original lemmas}
\Crefname{originallemma}{Original lemma}{Original lemmas}
\makeatletter
\newlabel{varieties-lemma-smooth-separable-closed-points-dense}{{056U}{source}{}{sourcepointer-4}{}}
\newlabel{varieties-lemma-smooth-separable-closed-points-dense@cref}{{[originallemma][0][]056U}{[1][0][]source}{}{}{}}
\makeatother
\title{Twisting to a one-dimensional space of sections}
\author{The Stacks Project Authors}
\date{stacks-project:pic.tex:local-snapshot-20261012; GNU FDL 1.2 or later}
\begin{document}
\maketitle
\noindent\textit{Editorial note.} Statement span is lines 1316-1326 of sources/stacks-pic/pinned-original-source.tex; proof span is lines 1328-1374, the complete author proof (45 body lines).\par\medskip
\noindent\textit{Editorial note.} Navigation stubs: 4 cited-but-not-included labels are supplied as non-load-bearing pointers with their pinned tags where the snapshot tag index has them.\par\medskip
\noindent\textit{Editorial note.} extra\_packages is empty on purpose: the retained preamble already loads xy with options and a second declaration clashes in the assembler wrapper.\par\medskip
\noindent\textit{Editorial note.} Text is retained verbatim, including the author's own typos and placeholders.\par\medskip
\noindent\textit{Editorial note.} Queue index 7 of queue\_pic.json (pb=2061, stmt 1316-1326, proof 1328-1374).\par\medskip
\noindent\textit{Editorial note.} stacks-pic wave-1 (pi01): fresh adjudication accept (2026-10-14); pool per standing queue hypercovering -> pic (SMC reassigned to Zhen per Chief 08:5x); pic.tex 27 labels/23 with proof judged 21 accept / 2 reject / 4 pb=0 env-dup rows excluded; spans cite-free and reference/slogan/history-free so no patchenv needed.\par\medskip
\noindent\textit{Allowed standard prerequisites (editorial).}\begin{itemize}
\item Explicit prerequisite boundary. The cited inputs are prerequisites; the argument in the retained span is not replaced by any cited result.
\end{itemize}\medskip
\ifdefined\hypertarget\hypertarget{sourcepointer-1}{}\fi
\noindent\textit{Source reference.} Original lemma 0B5U: sources/stacks-pic/pinned-original-source.tex, source lines 1374--1374. Cited by the retained statement or proof but not part of the collected spans; supplied only as a navigation pointer. The argument does not depend on its content.\par\medskip
\ifdefined\hypertarget\hypertarget{sourcepointer-2}{}\fi
\noindent\textit{Source reference.} Original lemma 0B8Z: sources/stacks-pic/pinned-original-source.tex, source lines 1374--1374. Cited by the retained statement or proof but not part of the collected spans; supplied only as a navigation pointer. The argument does not depend on its content.\par\medskip
\ifdefined\hypertarget\hypertarget{sourcepointer-3}{}\fi
\noindent\textit{Source reference.} Original lemma 0B8Y: sources/stacks-pic/pinned-original-source.tex, source lines 1374--1374. Cited by the retained statement or proof but not part of the collected spans; supplied only as a navigation pointer. The argument does not depend on its content.\par\medskip
\ifdefined\hypertarget\hypertarget{sourcepointer-4}{}\fi
\noindent\textit{Source reference.} Original lemma 056U: sources/stacks-pic/pinned-original-source.tex, source lines 1374--1374. Cited by the retained statement or proof but not part of the collected spans; supplied only as a navigation pointer. The argument does not depend on its content.\par\medskip
\begin{lemma}
\label{lemma-twist-with-general-divisor}
Let $k$ be a separably closed field. Let $X$ be a smooth projective
curve of genus $g$ over $k$. Let $K/k$ be a field extension and let
$\mathcal{L}$ be an invertible sheaf on $X_K$. Then there exists an
invertible sheaf $\mathcal{L}_0$ on $X$ such that
$\dim_K H^0(X_K,
\mathcal{L} \otimes_{\mathcal{O}_{X_K}} \mathcal{L}_0|_{X_K}) = 1$ and
$\dim_K H^1(X_K,
\mathcal{L} \otimes_{\mathcal{O}_{X_K}} \mathcal{L}_0|_{X_K}) = 0$.
\end{lemma}

\begin{proof}
This proof is a variant of the proof of
Varieties, Lemma \ref{varieties-lemma-general-degree-g-line-bundle}.
We encourage the reader to read that proof first.

\medskip\noindent
First we pick an ample invertible sheaf $\mathcal{L}_0$ and
we replace $\mathcal{L}$ by
$\mathcal{L} \otimes_{\mathcal{O}_{X_K}} \mathcal{L}_0^{\otimes n}|_{X_K}$
for some $n \gg 0$. The result will be that we may assume that
$H^0(X_K, \mathcal{L}) \not = 0$ and $H^1(X_K, \mathcal{L}) = 0$.
Namely, we will get the vanishing by Cohomology of Schemes, Lemma
\ref{coherent-lemma-vanshing-gives-ample} and the nonvanishing because
the degree of the tensor product is $\gg 0$.
We will finish the proof by descending induction on
$t = \dim_K H^0(X_K, \mathcal{L})$. The base case $t = 1$ is trivial.
Assume $t > 1$.

\medskip\noindent
Observe that for a $k$-rational point $x$ of $X$, the inverse image $x_K$
is a $K$-rational point of $X_K$. Moreover, there are infinitely many
$k$-rational points by Varieties, Lemma
\ref{varieties-lemma-smooth-separable-closed-points-dense}. Therefore
the points $x_K$ form a Zariski dense collection of points of $X_K$.

\medskip\noindent
Let $s \in H^0(X_K, \mathcal{L})$ be nonzero. From the previous paragraph
we deduce there exists a $k$-rational
point $x$ such that $s$ does not vanish in $x_K$. Let $\mathcal{I}$
be the ideal sheaf of $i : x_K \to X_K$ as in
Varieties, Lemma \ref{varieties-lemma-regular-point-on-curve}. Look at the
short exact sequence
$$
0 \to \mathcal{I} \otimes_{\mathcal{O}_{X_K}} \mathcal{L} \to
\mathcal{L} \to i_*i^*\mathcal{L} \to 0
$$
Observe that $H^0(X_K, i_*i^*\mathcal{L}) = H^0(x_K, i^*\mathcal{L})$
has dimension $1$ over $K$. Since $s$ does not vanish at $x$ we conclude that
$$
H^0(X_K, \mathcal{L}) \longrightarrow H^0(X, i_*i^*\mathcal{L})
$$
is surjective. Hence
$\dim_K H^0(X_K, \mathcal{I} \otimes_{\mathcal{O}_{X_K}} \mathcal{L}) = t - 1$.
Finally, the long exact sequence of cohomology also shows that
$H^1(X_K, \mathcal{I} \otimes_{\mathcal{O}_{X_K}} \mathcal{L}) = 0$
thereby finishing the proof of the induction step.
\end{proof}

\end{document}
